\section{Numerical Methods}
\label{sec:methods}

\subsection{Finite element discretisation}
All partial differential equations are solved using the open-source finite element library \texttt{FEniCSx} (dolfinx)~\cite{fenicsx2022}.
The computational domain is the rectangle \([0,2000]\times[0,500]\ \mu\text{m}\), discretised with an \(80\times20\) structured grid of triangular elements.
A coarse mesh of \(20\times5\) elements is used for rapid prototyping and proof-of-concept runs; mesh convergence is verified in the Supplementary Information (S1).

The Brinkman equations~\eqref{eq:brinkman} are discretised with the inf--sup stable Taylor--Hood element pair: continuous piecewise quadratic velocities (P2) and continuous piecewise linear pressure (P1).
This gives approximately \(20\,000\) degrees of freedom for the velocity block on the \(80\times20\) mesh.
The convection--diffusion equation~\eqref{eq:convdiff} is discretised with continuous piecewise linear elements (P1) and stabilised using the Streamline-Upwind Petrov--Galerkin (SUPG) method~\cite{brooks1982}.
The SUPG parameter \(\tau\) is evaluated element-wise using Eq.~\eqref{eq:supg}.
The design variable \(\rho\) and the filtered fields \(\tilde\rho,\bar\rho\) are also discretised with P1 elements on the same mesh.

The weak forms are assembled using the Unified Form Language (UFL) embedded in \texttt{FEniCSx}.
All volume integrals are evaluated with exact quadrature appropriate for the polynomial order; boundary integrals are evaluated over the tagged facets.

\subsection{Solver configuration}
During optimisation, the linear systems arising from the forward and adjoint Brinkman problems are solved with a direct LU factorisation (MUMPS) via PETSc’s \texttt{preonly}/\texttt{lu} combination.
This choice provides robust, parameter-free solutions for the moderate meshes used here (up to \(\sim 30\,000\) degrees of freedom).

For the scalability study presented in the Supplementary Information (S1), we additionally implement an iterative solver strategy: the Brinkman system is preconditioned with a block-triangular field-split preconditioner, using algebraic multigrid (HYPRE) for the velocity block and the pressure Schur complement~\cite{elman2014}.
The convection--diffusion system is preconditioned with HYPRE alone.
The scalability study confirms that the iterative approach maintains near-optimal convergence rates as the mesh is refined, enabling the pipeline to handle larger domains if needed.

\subsection{Optimisation algorithm}
The design is updated using the Optimality Criteria (OC) method~\cite{bendsoe2003}, which solves the Karush--Kuhn--Tucker conditions for a single volume constraint via bisection on the Lagrange multiplier \(\Lambda\).
The OC update formula reads
\begin{equation}
    \rho^{\text{new}} = \max\bigl(\rho_{\min},\,
                                 \min(\rho_{\max},\,
                                      \rho^{\text{old}} \cdot B^\eta)\bigr),
    \qquad B = -\frac{\mathrm{d}J/\mathrm{d}\rho}{\Lambda},
    \label{eq:oc}
\end{equation}
with damping exponent \(\eta = 0.5\) and a move limit of \(0.02\)--\(0.2\) to prevent divergence.
The Lagrange multiplier is determined by bisection until the volume constraint is satisfied to a tolerance of \(10^{-4}\).

For comparison, the Method of Moving Asymptotes (MMA)~\cite{svanberg1987} is also implemented via the \texttt{gcma} Python package; both update rules are provided in the code, and a convergence comparison is included in the Supplementary Information (S4).
The OC method was used for all results in the main paper.

\subsection{Continuation strategies}
To avoid premature trapping in poor local minima and to promote connected fluid pathways, we employ three nested continuation strategies.

\paragraph{Volume fraction continuation.}
The target fluid volume fraction \(V^*\) is linearly decreased from \(0.7\) to \(0.5\) over the first 40 iterations.
This initially relaxed constraint allows the optimiser to establish a well-connected channel network before removing excess fluid.

\paragraph{Sinusoidal initialisation.}
To break the symmetry of the uniform initial design---which would otherwise yield a vanishing sensitivity---the density field is initialised with a sinusoidal channel pattern:
\begin{equation}
  \rho(\mathbf{x})\big|_{t=0} = 0.7 + 0.25\sin\!\left(\frac{8\pi y}{L_y}\right),
  \label{eq:sinusoidal_init}
\end{equation}
clipped to \([0.01,\,0.99]\).
This creates four sinusoidal bands of alternating high and low density across the channel height, providing a non-zero sensitivity signal from the first iteration.

\paragraph{Concentration clamping.}
At high P\'{e}clet numbers, the SUPG-stabilised concentration field can still exhibit small non-physical oscillations (negative values or values slightly above 1) on coarse meshes.
After each forward solve, the concentration field is clipped to the physical range:
\begin{equation}
  c_h \leftarrow \max(0,\,\min(1,\,c_h)).
  \label{eq:clamping}
\end{equation}
This prevents oscillations from propagating into the adjoint right-hand side \(-\lambda\nabla c\), where they would otherwise cancel the sensitivity and stall convergence.

\paragraph{Heaviside sharpening.}
The projection parameter \(\beta\) (Eq.~\eqref{eq:heaviside}) is increased through the sequence
\[
\beta \in \{1,\,2,\,4,\,8,\,16,\,32,\,64\}
\]
every 20--40 iterations, gradually driving the design from grey to near-binary while avoiding the ill-conditioning that a sharp projection would cause early in the optimisation.

\paragraph{Brinkman penalty ramping.}
A third continuation is applied to the inverse permeability upper bound \(\alpha_{\max}\).
The value is ramped linearly from \(10^3\) to \(10^8\) over the first 60\% of the total iterations:
\begin{equation}
    \alpha_{\max} = 10^3 + (10^8 - 10^3)\cdot\min\!\left(1,\; \frac{\text{step}}{0.6\times \text{max\_iter}}\right).
    \label{eq:alpha_cont}
\end{equation}
This ramping ensures that fluid can flow through intermediate-density regions early in the optimisation, allowing the concentration field to develop and provide non-zero adjoint forcing.
\emph{Note:} in the \(80\times20\) runs reported here, \(\alpha_{\max}\) is held constant at \(10^{9}\,\text{Pa}\cdot\text{s}/\text{m}^2\) and symmetry breaking is instead achieved by the sinusoidal initialisation described above (Eq.~\eqref{eq:sinusoidal_init}).
The \(\alpha_{\max}\) ramping described in Eq.~\eqref{eq:alpha_cont} is recommended for finer meshes where the uniform-density starting point is more likely to stall.

\subsection{Post-processing and validation pipeline}
\label{sec:ns_validation}
After optimisation, the final physical density field \(\bar\rho\) is thresholded at \(\bar\rho = 0.5\) to obtain a binary fluid--solid distribution.
The fluid region is extracted using \texttt{pyvista} and remeshed with a body-fitted grid using \texttt{Gmsh}~\cite{geuzaine2009} to eliminate the Brinkman penalty approximation.
On this body-fitted mesh we solve the full incompressible Navier--Stokes equations (including the convective term) with a Picard iteration (five steps) to validate the design under real fluid physics.
The concentration field is then re-computed with SUPG-stabilised convection--diffusion on the same body-fitted mesh.
\paragraph{Outlet RMSE definition.}
\label{def:rmse}
The outlet concentration profile is compared to the prescribed target via the
\emph{root-mean-square error} (RMSE), defined as
%
\begin{equation}
  \mathrm{RMSE}
  \;=\;
  \sqrt{\frac{1}{N_{\mathrm{out}}}
        \sum_{i=1}^{N_{\mathrm{out}}}
        \bigl(c_h(y_i) - c_{\mathrm{target}}(y_i)\bigr)^{2}},
  \label{eq:rmse}
\end{equation}
%
where $N_{\mathrm{out}}$ is the number of P1 finite-element nodes on the outlet
boundary $\Gamma_{\mathrm{out}}$ ($x = L_x$), $c_h(y_i)$ is the computed
concentration at the $i$-th outlet node, and $c_{\mathrm{target}}(y_i)$ is the
prescribed target profile evaluated at the same node.

\emph{Normalisation.}
Both $c_h$ and $c_{\mathrm{target}}$ are dimensionless concentrations
normalised to the interval $[0,1]$, where $0$ corresponds to pure solvent
(inlet~1) and $1$ to pure solute (inlet~2).
RMSE is therefore dimensionless; a value of $0.098$ means that the
root-mean-square deviation from the target is $9.8$~percentage points of
the full $[0,1]$ concentration range.

\emph{Node sampling.}
For the $80\times20$ primary mesh the outlet boundary carries
$N_{\mathrm{out}} = n_y + 1 = 21$ uniformly spaced P1 nodes at
$y_i = (i-1)\,L_y/n_y$, $i=1,\ldots,21$, with spacing
$\Delta y = 25\,\mu\text{m}$.
For the $20\times5$ screening mesh, $N_{\mathrm{out}} = 6$
(spacing $\Delta y = 100\,\mu\text{m}$).
No additional interpolation is applied; RMSE is evaluated directly at the
mesh nodes produced by the P1 concentration solve.

\emph{Dimensional consistency.}
The RMSE formula~\eqref{eq:rmse} uses no dimensional quantities.
It is independent of the channel geometry and can be compared directly
across mesh resolutions provided the same normalisation convention is used.

Additional post-processing steps compute:
\begin{itemize}
    \item \textbf{Pressure drop} (fixed at \(p_{\text{in}}=1000\)~Pa) and total flow rate \(Q\);
    \item \textbf{Hydraulic resistance} \(R = \Delta p / Q\);
    \item \textbf{Maximum velocity} and \textbf{characteristic channel width} (via Euclidean distance transform);
    \item \textbf{Reynolds number} \(\text{Re} = \rho_f U_{\max} L_c / \mu\) and \textbf{Péclet number} \(\text{Pe} = U_{\max} L_c / D_{\text{fluid}}\);
    \item \textbf{Mixing length} -- the distance from the inlet where the concentration error drops below 5\% of the full scale;
    \item \textbf{Minimum feature size} -- the smallest channel width in the binary design, checked against a 10~\(\mu\text{m}\) photolithography limit.
\end{itemize}
All metrics are evaluated identically for the topology-optimised design and the classic Christmas tree network, the latter being generated as a density field on the same mesh using a set of line segments that replicate the geometry of Jeon \emph{et~al.}~\cite{jeon2000}.

\subsection{Reproducibility and open-source package}
\label{sec:reproducibility}

The complete pipeline is released as the open-source Python package
\texttt{micrograd} under the MIT license, available at
\url{https://github.com/nisongmonyimba278-byte/micrograd}.
The repository includes:

\begin{itemize}
    \item All finite element modules (mesh generation, forward and adjoint
          solvers, optimisation routines);
    \item Post-processing tools (STL export, body-fitted Navier--Stokes
          validation, experimental metrics);
    \item A Jupyter notebook for interactive use;
    \item A comprehensive test suite (validating the filter, forward solver
          against analytical Poiseuille flow, adjoint gradients via finite
          differences, MMA fallback logic, and mesh generation);
    \item A continuous-integration (CI) workflow
          (GitHub Actions~\cite{GitHubActions2024}) that automatically
          executes a three-stage pipeline on every commit:
          (i)~dependency installation inside the official
          \texttt{dolfinx/dolfinx:v0.7.3} Docker container, preserving
          the binary ABI compatibility required by the pre-compiled
          \textsc{PETSc} and \texttt{petsc4py} libraries by pinning
          \texttt{nlopt==2.7.1} and suppressing \textsc{NumPy} upgrades;
          (ii)~a forward-solver smoke test asserting $c_h\in[0,1]$ on a
          $20\times5$ mesh; and
          (iii)~a ten-iteration mini-optimisation verifying $J<10$,
          confirming end-to-end correctness of the adjoint sensitivity,
          OC update, and $\alpha$-continuation mechanism.
          During development, the CI pipeline identified and resolved
          three API incompatibilities introduced in \texttt{dolfinx~v0.7}:
          removal of the \texttt{petsc\_options\_prefix} keyword from
          \texttt{LinearProblem}, promotion of
          \texttt{element.interpolation\_points} from a property to a
          callable method, and a \textsc{NumPy}~2.0 binary incompatibility
          introduced by \texttt{nlopt}\,$\geq$\,2.8.
          All three were caught by the CI pipeline before affecting any
          scientific result.
          A passing CI badge on the repository landing page certifies
          that the complete solver stack is importable and produces
          physically meaningful output on a clean, containerised machine,
          independent of any local installation;
    \item An \texttt{environment.yaml} file pinning exact package versions
          to guarantee bit-identical reproduction.
\end{itemize}

The package is designed to run inside the official
\texttt{dolfinx/dolfinx:v0.7.3} Docker container, which provides a fully
self-contained environment with \mbox{FEniCSx~0.7.3}, \textsc{PETSc}, and
all necessary dependencies pre-compiled and version-locked.
The entire set of results presented in this paper can be reproduced with
a single command:

\begin{verbatim}
docker run --rm -v $(pwd):/root/micrograd \
  -e OMPI_MCA_btl=^openib \
  -e OMP_NUM_THREADS=1   \
  -e MPLBACKEND=Agg      \
  dolfinx/dolfinx:v0.7.3 bash -c \
  "cd /root/micrograd && bash run_all.sh"
\end{verbatim}

\noindent
The \texttt{run\_all.sh} script is environment-aware: it activates the
local Conda environment when available and falls back to the container's
system \texttt{python3} otherwise, so the same script drives both the
Docker workflow and the native \texttt{conda} workflow on Linux.

The code is permanently archived on Zenodo with
DOI~\texttt{10.5281/zenodo.XXXXXXX}~\cite{micrograd2025zenodo},
and a detailed API reference is hosted via GitHub Pages.


